
\documentclass[a4paper,12pt]{article}
\usepackage[T2A]{fontenc}
\usepackage[utf8]{inputenc}
\usepackage[russian]{babel}
\usepackage{amsmath,amssymb,mathtools}
\usepackage[scaled=0.92]{DejaVuSans}
\renewcommand{\familydefault}{\sfdefault}
\usepackage{icomma}
\usepackage[margin=20mm,top=18mm,bottom=19mm]{geometry}
\usepackage{microtype}
\usepackage{tikz}
\usetikzlibrary{calc,arrows.meta,positioning,shapes.geometric}
\usepackage{pgfplots}
\pgfplotsset{compat=1.18}
\usepackage{tabularx,booktabs,array}
\usepackage{enumitem}
\usepackage{hyperref}
\usepackage{fancyhdr}
\usepackage{setspace}
\usepackage{xcolor}

\definecolor{accent}{HTML}{087A7B}
\definecolor{darkaccent}{HTML}{124D55}
\definecolor{textgray}{HTML}{30383E}
\definecolor{muted}{HTML}{677780}
\definecolor{lightaccent}{HTML}{E4F0EF}
\color{textgray}
\hypersetup{hidelinks,pdftitle={Графики функций. Метод опорных точек},
pdfauthor={Артём Александрович Лёвин}}
\setstretch{1.05}
\setlength{\parindent}{0pt}
\setlength{\parskip}{3pt}
\setlist[itemize]{leftmargin=*,topsep=2pt,itemsep=1pt}
\setlist[enumerate]{leftmargin=*,topsep=2pt,itemsep=2pt}
\pagestyle{fancy}
\fancyhf{}
\fancyhead[L]{\small\color{muted} Графики функций · метод опорных точек}
\fancyhead[R]{\small\color{muted} 15.10.2026}
\fancyfoot[L]{\footnotesize\color{muted} Лёвин Артём Александрович}
\fancyfoot[R]{\footnotesize\color{muted} \thepage}
\renewcommand{\headrulewidth}{0.3pt}
\renewcommand{\footrulewidth}{0pt}
\setlength{\headheight}{15pt}
\newcommand{\sectiontitle}[1]{\vspace{4pt}{\Large\bfseries\color{darkaccent}#1}\par
 \vspace{1pt}{\color{accent}\rule{\linewidth}{0.55pt}}\vspace{4pt}}
\newcommand{\subsectiontitle}[1]{\vspace{3pt}{\large\bfseries\color{darkaccent}#1}\par\vspace{1pt}}
\newcommand{\key}[1]{{\color{darkaccent}\bfseries #1}}
\newcommand{\smallnote}[1]{{\small\color{muted}#1}}
\newcommand{\graphnote}[1]{{\small\color{muted}#1}}
\pgfplotsset{
  myaxis/.style={
    width=0.92\linewidth,
    height=5.7cm,
    axis lines=middle,
    axis line style={muted!90},
    tick style={muted!90},
    xlabel={$x$}, ylabel={$y$},
    xlabel style={font=\small}, ylabel style={font=\small},
    ticklabel style={font=\scriptsize},
    every axis plot/.append style={line width=1.2pt},
    grid=both, grid style={draw=lightaccent},
    minor grid style={draw=lightaccent!70},
    clip=true,
    samples=220,
    enlargelimits=false,
    legend style={draw=none,fill=none,font=\scriptsize,at={(0.98,0.98)},anchor=north east}
  }
}

\begin{document}
\hypersetup{pageanchor=false}
\begin{titlepage}
\thispagestyle{empty}
\begin{center}
\vspace*{22mm}
{\color{accent}\large МАТЕМАТИКА \textperiodcentered{} ЕГЭ ПРОФИЛЬНОГО УРОВНЯ}\par
\vspace{25mm}
{\Huge\bfseries\color{darkaccent} Чек-лист}\par
\vspace{7mm}
{\LARGE\bfseries Графики функций}\par
\vspace{2mm}
{\Large Метод опорных точек}\par
\vspace{13mm}
{\color{accent}\rule{79mm}{0.8pt}}\par
\vspace{13mm}
{\large Прочитайте точки графика.\\
Составьте формулу.\\
Ответьте точно на поставленный вопрос.}\par
\vspace{18mm}
{\large 15 октября 2026 года}\par
\vfill
{\normalsize\bfseries Лёвин Артём Александрович}\par
{\normalsize эксклюзивно для Ярослава Гаврилова}\par
\vspace{19mm}
\begin{tikzpicture}[x=1cm,y=1cm]
\draw[accent!65,line width=1.15pt] (-7,0) .. controls (-4,1.15) and (-2,-1) .. (0,0)
 .. controls (2,1) and (4,-1.15) .. (7,0);
\end{tikzpicture}
\end{center}
\end{titlepage}
\hypersetup{pageanchor=true}

\sectiontitle{01. Как извлечь формулу из графика}
\key{Главная идея.} Координаты точки $(x_0;y_0)$ графика функции
$y=f(x)$ удовлетворяют равенству $y_0=f(x_0)$.
Это позволяет найти неизвестные коэффициенты формулы.

\subsectiontitle{Как читать точки}
\begin{itemize}
\item \key{Сначала проверьте масштаб каждой оси.} Одна клетка необязательно равна единице. Особенно внимательно читайте ось $x$ у тригонометрических графиков: там бывают $\pi$, $\frac{\pi}{2}$, $\frac{\pi}{4}$.
\item Точку записывайте строго в порядке $(x;y)$: первая координата --- абсцисса, вторая --- ордината.
\item Выбирайте \key{опорные (узловые) точки}, координаты которых точно считываются по сетке. Они могут иметь дробные координаты и координаты с $\pi$.
\item Предпочитайте точки с $x=0$, $y=0$ или табличным значением аргумента: вычисления будут короче.
\end{itemize}

\subsectiontitle{Сколько точек нужно?}
Число искомых независимых коэффициентов подсказывает, сколько
независимых уравнений нужно составить. Для $y=a\cos x+b$
неизвестны $a,b$, поэтому обычно используют \key{две точки}.
Точки должны давать разные коэффициенты при $a$: например,
$\cos0=1$ и $\cos\pi=-1$.

\begin{center}
\renewcommand{\arraystretch}{1.24}
\begin{tabularx}{\linewidth}{@{}>{\raggedright\arraybackslash}p{37mm}
>{\raggedright\arraybackslash}p{62mm}X@{}}
\toprule
Вид графика & Общая формула & Коэффициенты\\
\midrule
Прямая & $y=kx+b$ & $k,\ b$\\
Косинусоида & $y=a\cos x+b$ & $a,\ b$\\
Синусоида & $y=a\sin x+b$ & $a,\ b$\\
Тангенсоида & $y=a\tg x+b$ & $a,\ b$\\
Парабола & $y=ax^2+bx+c$ & $a,\ b,\ c$\\
Гипербола со сдвигом & $y=\dfrac{a}{x-h}+k$ & $a,\ h,\ k$\\
\bottomrule
\end{tabularx}
\end{center}
\smallnote{В последних двух строках приведены общие формы для распознавания.
Для смещённой гиперболы дополнительно требуется $x\ne h$;
её коэффициенты не всегда удобно искать простой линейной системой.}

\subsectiontitle{Обязательная логика решения}
\begin{enumerate}
\item Узнайте семейство функции и выпишите его общую формулу.
\item Выберите точки с удобными, точно считываемыми координатами.
\item Для каждой точки запишите равенство $y_i=f(x_i)$.
\item Вычислите известные значения функций и решите систему.
\item Подставьте найденные коэффициенты в общую формулу.
\item Прочитайте вопрос ещё раз: требуется коэффициент, $f(x_0)$,
неизвестный $x$ или точка пересечения?
\end{enumerate}
\clearpage
\thispagestyle{fancy}
\vspace*{5mm}
\begin{center}
{\LARGE\bfseries\color{darkaccent} Алгоритм: метод опорных точек}\par
\vspace{3mm}
{\small\color{muted} Универсальный маршрут для задач на восстановление формулы}
\end{center}
\vspace{6mm}
\tikzset{
flowstart/.style={ellipse,draw=accent,fill=lightaccent,
 line width=.85pt,minimum width=57mm,minimum height=11mm,align=center,font=\small\bfseries},
flowact/.style={rounded corners=3pt,draw=accent,line width=.8pt,
 text width=100mm,minimum height=12mm,align=center,inner sep=4pt,font=\small},
flowcheck/.style={diamond,aspect=2.85,draw=accent,line width=.8pt,
 fill=lightaccent,text width=59mm,align=center,inner sep=1.8pt,font=\small\bfseries},
flowarrow/.style={-{Stealth[length=2.4mm]},draw=darkaccent,line width=.8pt},
flowlabel/.style={font=\scriptsize\bfseries,color=darkaccent}
}
\begin{center}
\begin{tikzpicture}[x=1cm,y=1cm]
\node[flowstart] (a) at (0,0) {Прочитайте условие};
\node[flowact] (b) at (0,-1.95) {Определите вид функции и неизвестные коэффициенты};
\node[flowact] (c) at (0,-3.90) {Проверьте масштаб осей; выберите удобные точки};
\node[flowact] (d) at (0,-5.85) {Для каждой точки запишите $y_i=f(x_i)$};
\node[flowact] (e) at (0,-7.80) {Вычислите значения и решите систему};
\node[flowcheck] (f) at (0,-10.0) {Какой результат требуется?};
\node[flowact,text width=49mm,minimum height=18mm] (left) at (-3.55,-12.75)
 {Коэффициент?\\Запишите его значение};
\node[flowact,text width=49mm,minimum height=18mm] (right) at (3.55,-12.75)
 {Функция, аргумент или\\пересечение?\\Найдите формулы};
\node[flowact,text width=90mm,minimum height=14mm] (g) at (0,-15.35)
 {Подставьте значение или приравняйте функции};
\node[flowstart] (h) at (0,-18.15) {Сверьте результат с вопросом};
\draw[flowarrow](a)--(b);
\draw[flowarrow](b)--(c);
\draw[flowarrow](c)--(d);
\draw[flowarrow](d)--(e);
\draw[flowarrow](e)--(f);
\draw[flowarrow](f.south west)--node[left,flowlabel]{коэффициент}(left.north);
\draw[flowarrow](f.south east)--node[right,flowlabel]{другой результат}(right.north);
\draw[flowarrow](left.south) -- ++(0,-.55) -| (-5.2,-16.7) |- (h.west);
\draw[flowarrow](right.south)--(g.north east);
\draw[flowarrow](g.south)--(h.north);
\end{tikzpicture}
\end{center}

\clearpage
\sectiontitle{02. График косинуса: два опорных значения}
\subsectiontitle{Разобранный пример: $y=a\cos x+b$}
По графику считаны точки $A(0;\frac12)$ и $B(\pi;-\frac52)$.
Подставим их координаты:
\[
\begin{cases}
\frac12=a\cos0+b,\\[2pt]
-\frac52=a\cos\pi+b.
\end{cases}
\quad\Longleftrightarrow\quad
\begin{cases}
a+b=\frac12,\\[2pt]
-a+b=-\frac52.
\end{cases}
\]
Вычитаем второе уравнение из первого:
\[
2a=\frac12+\frac52=3,\qquad
\boxed{a=\frac32}.
\]
Затем $b=\frac12-\frac32=-1$. Полная формула:
\[
\boxed{f(x)=\frac32\cos x-1}.
\]
\begin{center}
\begin{tikzpicture}
\begin{axis}[
myaxis,
 height=7.0cm,
 xmin=-1.8,xmax=5.0,
 ymin=-3.2,ymax=0.8,
 xtick={-1.5708,0,1.5708,3.1416,4.7124},
 xticklabels={$-\frac{\pi}{2}$,$0$,$\frac{\pi}{2}$,$\pi$,$\frac{3\pi}{2}$},
 ytick={-3,-2,-1,0},
 legend entries={$y=\frac32\cos x-1$}
]
\addplot[accent,domain=-1.57:4.71]{1.5*cos(deg(x))-1};
\addplot[only marks,mark=*,mark size=2.2pt,darkaccent] coordinates {(0,0.5) (3.1416,-2.5)};
\node[font=\scriptsize,anchor=south west] at (axis cs:0,0.5) {$A(0;\frac12)$};
\node[font=\scriptsize,anchor=north west] at (axis cs:3.1416,-2.5) {$B(\pi;-\frac52)$};
\end{axis}
\end{tikzpicture}
\end{center}
\graphnote{На графике видно, почему удобно брать точки с $x=0$ и $x=\pi$: значения косинуса табличные, а координаты читаются без двусмысленности.}

\subsectiontitle{Случай, когда неизвестный исчезает}
Пусть график $y=a\cos x+b$ проходит через
$A(0;\frac32)$ и $B(\frac\pi2;-\frac12)$.
Тогда
\[
\begin{cases}
\frac32=a\cos0+b,\\
-\frac12=a\cos\frac\pi2+b
\end{cases}
\Longleftrightarrow
\begin{cases}
\frac32=a+b,\\
-\frac12=b.
\end{cases}
\]
Следовательно, $\boxed{b=-\frac12}$ и $\boxed{a=2}$.
Здесь полезнее начинать со второго уравнения:
$\cos\frac\pi2=0$, поэтому коэффициент $a$ исчезает.

\clearpage
\sectiontitle{03. График синуса: знаки и периодичность}
\subsectiontitle{Разобранный пример: $y=a\sin x+b$}
Пусть известны $A(0;-1)$ и $B(-\frac\pi2;1)$.
Тогда
\[
\begin{cases}
-1=a\sin0+b,\\
1=a\sin(-\frac\pi2)+b
\end{cases}
\Longleftrightarrow
\begin{cases}
-1=b,\\
1=-a+b.
\end{cases}
\]
Отсюда $b=-1$, $1=-a-1$, значит $a=-2$:
\[
\boxed{f(x)=-2\sin x-1}.
\]
При необходимости вычислить $f(\frac{25\pi}{6})$:
\[
\frac{25\pi}{6}=4\pi+\frac\pi6,\quad
\sin\frac{25\pi}{6}=\sin\frac\pi6=\frac12,\quad
f\!\left(\frac{25\pi}{6}\right)=\boxed{-2}.
\]
\begin{center}
\begin{tikzpicture}
\begin{axis}[
myaxis,
 height=8.3cm,
 xmin=-3.4,xmax=6.6,
 ymin=-3.4,ymax=1.4,
 xtick={-3.1416,-1.5708,0,1.5708,3.1416,4.7124,6.2832},
 xticklabels={$-\pi$,$-\frac{\pi}{2}$,$0$,$\frac{\pi}{2}$,$\pi$,$\frac{3\pi}{2}$,$2\pi$},
 ytick={-3,-2,-1,0,1},
 legend entries={$y=-2\sin x-1$}
]
\addplot[accent,domain=-3.1416:6.2832]{-2*sin(deg(x))-1};
\addplot[only marks,mark=*,mark size=2.2pt,darkaccent] coordinates {(0,-1) (-1.5708,1)};
\node[font=\scriptsize,anchor=south west] at (axis cs:0,-1) {$A(0;-1)$};
\node[font=\scriptsize,anchor=south east] at (axis cs:-1.5708,1) {$B(-\frac{\pi}{2};1)$};
\end{axis}
\end{tikzpicture}
\end{center}
\graphnote{График помогает проконтролировать знак коэффициента $a$: при переходе от $x=-\frac\pi2$ к $x=0$ функция убывает, значит перед синусом действительно стоит отрицательный множитель.}

\clearpage
\sectiontitle{04. Тангенс: когда ответ находится сразу}
\subsectiontitle{Разобранный пример: $y=a\tg x+b$}
График проходит через $A(0;-\frac32)$ и
$B(\frac\pi4;\frac12)$:
\[
\begin{cases}
-\frac32=a\tg0+b,\\
\frac12=a\tg\frac\pi4+b
\end{cases}
\Longleftrightarrow
\begin{cases}
-\frac32=b,\\
\frac12=a+b.
\end{cases}
\]
Если требуется найти $b$, ответ \key{уже получен}:
$\boxed{b=-\frac32}$. Если нужна вся формула, то
\[
a=\frac12+\frac32=2,\qquad
\boxed{f(x)=2\tg x-\frac32}.
\]
\smallnote{Область определения тангенса:
$x\ne\frac\pi2+\pi n$, где $n\in\mathbb Z$.}
\begin{center}
\begin{tikzpicture}
\begin{axis}[
myaxis,
 height=8.6cm,
 xmin=-1.8,xmax=1.8,
 ymin=-6.2,ymax=3.2,
 xtick={-1.5708,0,0.7854,1.5708},
 xticklabels={$-\frac{\pi}{2}$,$0$,$\frac{\pi}{4}$,$\frac{\pi}{2}$},
 ytick={-6,-4,-2,0,2},
 legend entries={$y=2\tg x-\frac32$}
]
\addplot[muted,dashed] coordinates {(-1.5708,-6.2) (-1.5708,3.2)};
\addplot[muted,dashed] coordinates {(1.5708,-6.2) (1.5708,3.2)};
\addplot[accent,domain=-1.45:1.45]{2*tan(deg(x))-1.5};
\addplot[only marks,mark=*,mark size=2.2pt,darkaccent] coordinates {(0,-1.5) (0.7854,0.5)};
\node[font=\scriptsize,anchor=north west] at (axis cs:0,-1.5) {$A(0;-\frac32)$};
\node[font=\scriptsize,anchor=south west] at (axis cs:0.7854,0.5) {$B(\frac{\pi}{4};\frac12)$};
\end{axis}
\end{tikzpicture}
\end{center}
\graphnote{Пунктиром показаны вертикальные асимптоты. На таком графике особенно важно брать точки внутри области определения.}

\clearpage
\sectiontitle{05. Прямая: та же система без тригонометрии}
\subsectiontitle{Разобранный пример: $y=kx+b$}
По графику считываем $A(-1;-3)$ и $B(3;4)$:
\[
\begin{cases}
-3=-k+b,\\
4=3k+b.
\end{cases}
\]
Вычтем первое равенство из второго:
\[
7=4k,\quad k=\frac74,\qquad
b=-3+\frac74=-\frac54.
\]
Формула восстановлена:
\[
\boxed{f(x)=\frac74x-\frac54}.
\]
\begin{center}
\begin{tikzpicture}
\begin{axis}[
myaxis,
 height=8.0cm,
 xmin=-3.4,xmax=4.4,
 ymin=-5.2,ymax=5.4,
 xtick={-3,-2,-1,0,1,2,3,4},
 ytick={-5,-4,-3,-2,-1,0,1,2,3,4,5},
 legend entries={$y=\frac74x-\frac54$}
]
\addplot[accent,domain=-3.2:4.2]{1.75*x-1.25};
\addplot[only marks,mark=*,mark size=2.2pt,darkaccent] coordinates {(-1,-3) (3,4)};
\node[font=\scriptsize,anchor=south west] at (axis cs:-1,-3) {$A(-1;-3)$};
\node[font=\scriptsize,anchor=south east] at (axis cs:3,4) {$B(3;4)$};
\end{axis}
\end{tikzpicture}
\end{center}

\subsectiontitle{Два разных вопроса к одному графику}
\begin{itemize}
\item \key{«Найдите $f(-5)$».} Известен аргумент:
\[
f(-5)=\frac74\cdot(-5)-\frac54
=-\frac{40}{4}=\boxed{-10}.
\]
\item \key{«Найдите $x$, при котором $f(x)=-13{,}5$».}
Известно значение функции; требуется решить уравнение:
\[
-\frac{27}{2}=\frac74x-\frac54
\Longleftrightarrow -54=7x-5
\Longleftrightarrow \boxed{x=-7}.
\]
\end{itemize}
\key{Запомните различие:}
«найдите $f(x_0)$» означает подстановку $x_0$;
«найдите $x$ при $f(x)=c$» означает решение уравнения.

\clearpage
\sectiontitle{06. Пересечение двух графиков}
Если показаны две прямые, формулы которых в условии отсутствуют,
каждую можно восстановить по двум точкам:
\[
f_1(x)=k_1x+b_1,\qquad f_2(x)=k_2x+b_2.
\]
В общей точке графиков их ординаты совпадают:
\[
\boxed{f_1(x)=f_2(x)}.
\]
\subsectiontitle{Разобранный пример}
Пусть после восстановления формул получено
\[
f_1(x)=4x+9,\qquad
f_2(x)=\frac32x-\frac72.
\]
Приравниваем правые части:
\[
4x+9=\frac32x-\frac72
\quad\Longleftrightarrow\quad
\frac52x=-\frac{25}{2}
\quad\Longleftrightarrow\quad
\boxed{x=-5}.
\]
Это \key{абсцисса} точки пересечения.
Если нужна ещё и ордината, подставляем $x=-5$:
\[
y=f_1(-5)=4(-5)+9=\boxed{-11}.
\]
Точка пересечения: $\boxed{(-5;-11)}$.
Приблизительное положение точки на рисунке полезно для проверки знака и порядка величины.
\begin{center}
\begin{tikzpicture}
\begin{axis}[
myaxis,
 height=8.8cm,
 xmin=-7.2,xmax=1.2,
 ymin=-15.2,ymax=11.8,
 xtick={-7,-6,-5,-4,-3,-2,-1,0,1},
 ytick={-15,-10,-5,0,5,10},
 legend entries={$y=4x+9$,$y=\frac32x-\frac72$}
]
\addplot[accent,domain=-7.2:1.0]{4*x+9};
\addplot[darkaccent,domain=-7.2:1.0]{1.5*x-3.5};
\addplot[only marks,mark=*,mark size=2.5pt,black] coordinates {(-5,-11)};
\node[font=\scriptsize,anchor=north west] at (axis cs:-5,-11) {$(-5;-11)$};
\end{axis}
\end{tikzpicture}
\end{center}
\graphnote{На рисунке точка пересечения находится левее нуля и ниже оси $x$. Такой визуальный контроль помогает заметить грубую вычислительную ошибку.}

\clearpage
\sectiontitle{07. Тригонометрический минимум для графиков}
\begin{center}
\renewcommand{\arraystretch}{1.30}
\begin{tabular}{@{}lccccc@{}}
\toprule
$x$ & $-\frac\pi2$ & $0$ & $\frac\pi4$ & $\frac\pi2$ & $\pi$\\
\midrule
$\sin x$ & $-1$ & $0$ & $\frac{\sqrt2}{2}$ & $1$ & $0$\\
$\cos x$ & $0$ & $1$ & $\frac{\sqrt2}{2}$ & $0$ & $-1$\\
$\tg x$ & --- & $0$ & $1$ & --- & $0$\\
\bottomrule
\end{tabular}
\end{center}
\[
\sin(x+2\pi n)=\sin x,\qquad
\cos(x+2\pi n)=\cos x,\qquad
\tg(x+\pi n)=\tg x.
\]
\key{Чётные и нечётные кратные $\pi$:}
$\cos(2\pi n)=1$; $\cos((2n+1)\pi)=-1$;
синус любого целого кратного $\pi$ равен нулю.

\subsectiontitle{Как вычислять сложный аргумент}
Пусть $f(x)=2\cos x-\frac12$, требуется $f(\frac{13\pi}{4})$.
Сначала считаем косинус:
\[
\cos\frac{13\pi}{4}
=\cos\!\left(3\pi+\frac\pi4\right)
=-\frac{\sqrt2}{2}.
\]
Поэтому
\[
f\!\left(\frac{13\pi}{4}\right)
=2\!\left(-\frac{\sqrt2}{2}\right)-\frac12
=\boxed{-\sqrt2-\frac12}.
\]

\sectiontitle{08. Что важно удержать из прикладных задач}
\subsectiontitle{Дифракционная решётка}
Для максимума порядка $m$:
\[
d\sin\varphi=m\lambda.
\]
Здесь $d$ --- период решётки, $\lambda$ --- длина волны,
$m$ --- номер максимума, $\varphi$ --- угол.
Например, при $d=2000$ нм, $\lambda=500$ нм, $m=2$:
\[
\sin\varphi=\frac{2\cdot500}{2000}=\frac12
\quad\Longrightarrow\quad
\boxed{\varphi=30^\circ}
\quad (0^\circ\le\varphi\le90^\circ).
\]
Условие существования максимума: $m\lambda\le d$.

\subsectiontitle{Квадрат синуса и выбор угла}
Если прикладная формула имеет вид
$Q=mv^2\sin^2\alpha$, то при $Q=75$, $m=3$, $v=10$
\[
\sin^2\alpha=\frac{75}{3\cdot100}
=\frac14.
\]
На интервале $0^\circ\le\alpha\le90^\circ$
$\sin\alpha\ge0$, поэтому
\[
\sin\alpha=\frac12,\qquad \boxed{\alpha=30^\circ}.
\]
\key{Важно:} из $\sin^2\alpha=\frac14$ без дополнительных
условий следует $|\sin\alpha|=\frac12$.
Знак и окончательный угол определяются заданным интервалом.

\subsectiontitle{Корень в знаменателе}
Приём рационализации знаменателя:
\[
\frac{1}{\sqrt2}
=\frac{1}{\sqrt2}\cdot\frac{\sqrt2}{\sqrt2}
=\boxed{\frac{\sqrt2}{2}}.
\]
Именно это табличное значение синуса и косинуса угла $45^\circ$.

\sectiontitle{09. Финальная самопроверка}
\subsectiontitle{Какую операцию выбирать по формулировке?}
\begin{center}
\renewcommand{\arraystretch}{1.24}
\begin{tabularx}{\linewidth}{@{}>{\bfseries\raggedright\arraybackslash}p{54mm}X@{}}
\toprule
Формулировка & Что делать\\
\midrule
«Найдите коэффициент» & Остановиться, как только коэффициент найден из системы.\\
«Найдите $f(x_0)$» & Сначала восстановить формулу, затем подставить известный $x_0$.\\
«Найдите $x$, при котором $f(x)=c$» & Восстановить формулу и решить уравнение относительно $x$.\\
«Найдите абсциссу точки пересечения» & Восстановить формулы обеих функций и приравнять их.\\
«Найдите ординату точки пересечения» & После нахождения $x$ подставить его в любую из функций.\\
\bottomrule
\end{tabularx}
\end{center}

\subsectiontitle{Контроль перед записью ответа}
\begin{itemize}
\item Верно ли считан масштаб обеих осей, включая обозначения с $\pi$?
\item Правильно ли записаны координаты в порядке $(x;y)$?
\item Дают ли опорные точки независимые уравнения?
\item Не перепутаны ли $\sin0=0$ и $\cos0=1$?
\item Не потерян ли минус при $\cos\pi=-1$ или $\sin(-\frac\pi2)=-1$?
\item Доведена ли система лишь до той величины, которую просят найти?
\item Различены ли $f(x_0)$, неизвестный $x$ и абсцисса пересечения?
\item Проверены ли ограничения области определения и диапазон углов?
\end{itemize}
\vfill
{\color{accent}\rule{\linewidth}{0.5pt}}\par
{\small\color{muted} Лёвин Артём Александрович · эксклюзивно для Ярослава Гаврилова · 15.10.2026}
\end{document}
